% Generated draft: requires mathematical and source review.
\documentclass[11pt]{article}
\usepackage[margin=25mm]{geometry}
\usepackage{fontspec}
\setmainfont{Latin Modern Roman}
\usepackage{amsmath,amssymb,mathtools,unicode-math}
\setmathfont{Latin Modern Math}
\usepackage{enumitem,needspace,xurl,hyperref}
\setlength{\parindent}{0pt}
\setlength{\parskip}{5pt}
\setlength{\emergencystretch}{3em}
\newcommand{\N}{\mathbb{N}}
\newcommand{\Z}{\mathbb{Z}}
\newcommand{\Q}{\mathbb{Q}}
\newcommand{\R}{\mathbb{R}}
\newcommand{\C}{\mathbb{C}}
\newcommand{\F}{\mathbb{F}}
\newcommand{\A}{\mathbb{A}}
\newcommand{\PP}{\mathbb{P}}
\newcommand{\E}{\mathbb{E}}
\newcommand{\eps}{\varepsilon}
\newcommand{\dd}{\,\mathrm{d}}
\begin{document}
\section*{UnsolvedMath ID: 20002313}
\textbf{Source problem code:} catalogue entry
\subsection{Definitions and mathematical statement}
The source question asks: let $A=HF(\C)$, and let $K\cong\C$ be a $\Sigma$-definable field in $A$. Must there be a $\Sigma$-definable field isomorphism $\C\to K$? Here $HF(\C)$ denotes a hereditarily finite superstructure with complex urelements and the complex field structure. The exact language, meaning of $\Sigma$, permitted parameters, and whether a definable field has a literal domain or an interpreted quotient must be fixed from the source before this becomes a fully specified formal question; they are not supplied in this record.

The accompanying comparison says that the answer is positive with $\R$ in place of $\C$ when the field copy has domain $K\subseteq\R$, whereas the case $K\subseteq\R^2$ is reported open. The extracted real-field display loses a binary-operation symbol; its intended reading as a field copy $(K,\oplus,\otimes)\cong(\R,+,\cdot)$ is a reconstruction, not a verified transcription here. These assertions have the scope of the workshop account, rather than an independently checked current literature assessment.

The research aid supplies the following obstruction, which concerns literal domains and ordinary definability with finitely many parameters. For $a\in HF(\C)$, let $\operatorname{supp}(a)$ be the finite set of complex urelements occurring in its finite membership closure. Suppose field automorphisms of $\C$ act recursively on $HF(\C)$ as automorphisms of the chosen language. Let $\bar p$ be a finite parameter tuple and put $S=\bigcup_i\operatorname{supp}(p_i)$. If $f:\C\to D\subseteq HF(\C)$ is a $\bar p$-definable bijection, then
\[
f(\sigma z)=\sigma f(z)\qquad(\sigma\in\operatorname{Aut}(\C/S)).
\]
For $T_z=\operatorname{supp}(f(z))$, the aid asserts
\[
T_z\subseteq\operatorname{acl}_{\C}(S\cup\{z\}),\qquad
z\in\operatorname{acl}_{\C}(S\cup T_z),\qquad
\operatorname{trdeg}_{\Q(S)}\Q(S,T_z)\leq1.
\]
Here $\operatorname{acl}_{\C}(U)$ means the elements algebraic over the subfield $\Q(U)$. Thus, if for every finite $S\subseteq\C$ the literal domain $D$ contains $d$ with $\operatorname{trdeg}_{\Q(S)}\Q(S,\operatorname{supp}(d))\geq2$, no such finitely parameter definable bijection, and consequently no such definable field isomorphism, exists. This is a conditional obstruction; the supplied record does not construct a $\Sigma$-definable field copy satisfying its hypothesis. It does not settle the original question or automatically apply to quotient interpretations.

The trailing text ``1.6. $\lambda$-many models of each cardinality $\lambda\geq\aleph_1$ (Greenberg)'' belongs to the next section heading and is not an additional hypothesis of Question 11.
\subsection{Short English statement}
Question 11 asks whether a $\Sigma$-definable copy of the complex field inside $HF(\C)$ admits a $\Sigma$-definable isomorphism with $\C$. The supplied aid gives a finite-support obstruction for literal domains; it supplies neither a counterexample satisfying that obstruction nor a resolution of the question. The real-field comparison and extraction defects require source review.
\subsection{Sources}
\begin{itemize}
\item Local UnsolvedMath registry, permanent ID $20002313$, source code \texttt{catalogue entry}; canonical file \texttt{aim-logic-notes.json}, \texttt{notes[88]}.
\item American Institute of Mathematics, \emph{Computable stability theory}, Question 11, \url{https://aimath.org/pastworkshops/computestabproblems.pdf}. No proposer or full bibliographic reference is supplied. The verification field reports inspection of page 3 (PDF page index 2) and its math-font encoding, but does not actually include the promised recovered passage.
\item Supplied research attempt 1: partial result, problem status unknown at run, rigor label ``proved,'' novelty confidence low. Registry classification: partial progress; registry status: partially solved; statement status: reconstructed, unverified. The finite-support obstruction is attributed only to this machine-generated research aid.
\end{itemize}
\end{document}
